% Hafta 12 — Sızma testi planının dört başlığı
% Derleme: py -3.12 tools/latex/derle.py docs/week-12/assets/h12-11-pentest-plani.tex
\documentclass[tikz,border=8pt]{standalone}
\usepackage{cen429-cizim}
\begin{document}
\begin{tikzpicture}[node distance=10mm and 10mm]
  \node[baslik] (b) at (0,0) {Sızma testi planının dört başlığı};
  \node[altbaslik, text width=160mm] at (b.south west) {plan olmadan yapılan test ne tekrarlanabilir ne de savunulabilir};

  \node[kutu=32mm, below=16mm of b.south west, anchor=north west] (p1) {\kb{1 · KAPSAM}\\ne test edilir\\ne edilmez};
  \node[kotukutu=38mm, right=of p1] (p2) {\kb{2 · ANGAJMAN KURALLARI}\\izin, sınır, zaman\\DURDURMA KOŞULU};
  \node[kutu=30mm, right=of p2] (p3) {\kb{3 · YÖNTEM}\\WSTG / PTES / MASTG};
  \node[iyikutu=40mm, right=of p3] (p4) {\kb{4 · TEST KARTLARI}\\tekrarlanabilir\\adım + beklenen +\\gözlenen};

  \draw[ok, draw=soluk] (p1) -- (p2);
  \draw[ok, draw=soluk] (p2) -- (p3);
  \draw[ok, draw=soluk] (p3) -- (p4);

  \node[dip, text width=160mm, below=8mm of p4.south -| p2, anchor=north] {%
    Durdurma koşulu, gerçek zarar/veri kaybı riskini sınırlar --- etik ve güvenli yürütme.};
\end{tikzpicture}
\end{document}
